The magic functions $a, b : \R^8 \to \C$ are radial: they are built from profiles
$f : \R \to \C$ via $x \mapsto f(\|x\|^2)$. This section describes the bridge that converts
smoothness and one-sided decay estimates of a profile $f$ on $[0, \infty)$ into the statement
that $x \mapsto f(\|x\|^2)$ is a Schwartz function on $\R^8$.

The profiles arising from the contour integrals are only controlled for $r \geq 0$ (which
suffices, since $\|x\|^2 \geq 0$); to obtain a Schwartz function on all of $\R$ we first multiply
by a smooth cutoff supported near $[0, \infty)$.

\begin{definition}[Smooth transition function]\label{def:smooth-transition}\leanok
  A \emph{smooth transition function} is a $C^\infty$ function $\chi : \R \to \R$ such that
  \[
    \chi(r) = 0 \text{ for } r \le -\tfrac12, \qquad
    \chi(r) = 1 \text{ for } r \ge 0,
  \]
  and $0 \le \chi(r) \le 1$ for all $r \in \R$.
\end{definition}

\begin{remark}[Mathlib realization]\label{def:radial-cutoff}
  In Mathlib, the function $\chi$ is explicitly constructed as $\chi(r) := \operatorname{smoothTransition}(2r + 1)$,
  where $\operatorname{smoothTransition}$ is the standard $C^\infty$ transition function on $\R$.
  We denote by $\chi_\C : \R \to \C$ the complex-valued version $\chi_\C(r) := \chi(r)$ (viewed as a
  complex number). This is the definition \lean{RadialSchwartz.cutoffC}\leanok in the formalization.
\end{remark}

\begin{lemma}[Cutoff profiles are globally Schwartz]\label{lemma:cutoff-mul-decay}\uses{def:smooth-transition}\lean{RadialSchwartz.cutoffC_mul_decay_of_nonneg_of_contDiff, RadialSchwartz.contDiff_cutoffC_mul_of_contDiffOn_Ioi_neg1}\leanok
  Let $f : \R \to \C$.
  \begin{enumerate}
    \item If $f$ is $C^\infty$ on $(-1, \infty)$, then $\chi f$ is $C^\infty$ on $\R$.
    \item If moreover $\chi f$ is smooth and $f$ satisfies the one-sided Schwartz bounds
      \[
        \forall k, n \in \N, \ \exists C, \ \forall x \ge 0, \quad
          |x|^k \, \bigl\| (f)^{(n)}(x) \bigr\| \le C,
      \]
      then $\chi f$ satisfies the same bounds for \emph{all} $x \in \R$.
  \end{enumerate}
\end{lemma}
\begin{proof}\leanok
  For (1): near any $x < -\tfrac12$ the product vanishes identically; elsewhere both factors are
  smooth. For (2): for $x \le -1$ all iterated derivatives vanish; on the compact interval
  $[-1, 0]$ the continuous function $|x|^k \|(\chi f)^{(n)}(x)\|$ is bounded; for $x > 0$ the
  product agrees with $f$ near $x$ and the one-sided bounds apply.
\end{proof}

\begin{definition}[Schwartz packaging of a profile]\label{def:fCutSchwartz}\uses{def:smooth-transition, lemma:cutoff-mul-decay, def:Schwartz-Space}\lean{RadialSchwartz.Bridge.fCutSchwartz}\leanok
  Given $f : \R \to \C$ smooth with one-sided Schwartz bounds on $[0, \infty)$, the function
  $\chi f$ is a Schwartz function on $\R$ by Lemma~\ref{lemma:cutoff-mul-decay}; we denote the
  resulting element of $\mathcal{S}(\R, \C)$ by $f^{\mathrm{cut}}$. It agrees with $f$ on
  $[0, \infty)$.
\end{definition}

\begin{theorem}[Radial lift to $\R^d$]\label{thm:schwartz-multidimensional}\uses{def:Schwartz-Space}\lean{schwartzMap_multidimensional_of_schwartzMap_real}\leanok
  Let $F$ be a real inner product space (for us, $F = \R^8$) and let
  $g \in \mathcal{S}(\R, \C)$. Then
  \[
    x \mapsto g(\|x\|^2)
  \]
  is a Schwartz function on $F$.
\end{theorem}
\begin{proof}\leanok
  The map $x \mapsto \|x\|^2$ has temperate growth, and composition of a Schwartz function with a
  temperate map along a suitable polynomial bound preserves the Schwartz property (Mathlib's
  \texttt{SchwartzMap.compCLM}).
\end{proof}

Combining Definition~\ref{def:fCutSchwartz} with Theorem~\ref{thm:schwartz-multidimensional}: if
a radial profile $f : \R \to \C$ is smooth and satisfies Schwartz-type bounds for all iterated
derivatives on $r \ge 0$, then $x \mapsto f(\|x\|^2)$ is a Schwartz function on $\R^8$. This is
how the functions $a$ and $b$ are assembled from the profiles $I_1', \dots, I_6'$ and
$J_1', \dots, J_6'$ in the next section.
